% Comando simples para exibir comandos Latex no texto
\newcommand{\comando}[1]{\textbf{$\backslash$#1}}
 

%Atualmente, uma grande discussão que existe relacionado a logística é sobre a localização de facilidades, sendo esta responsável pelo dimensionamento da quantidade e localização geográfica de facilidades, para que suportem uma cadeia de suprimentos no que diz respeito a manutenção de itens em estoque, transporte e operação.

Os problemas de localização de facilidades ou \sigla{FLP}{\emph{Facility Location Problem}} são de extrema importância estratégica para uma empresa, pois com um dimensionamento adequado é possível ter um atendimento mais ágil e com menor custo, tendo vantagem competitiva considerável, pois os maiores custos de uma empresa são os custos logísticos \cite{Owen1998Strategic}.

Esse tipo de problema não se restringe apenas a aplicações voltadas para a redução de custos no setor logístico; ele possui diversas aplicações em diferentes áreas e com objetivos variados. Na área da saúde, por exemplo, pode ser usado para ampliar a cobertura de atendimento em uma população. Em regiões propensas a desastres naturais, o foco é criar áreas seguras que possam atender à população rapidamente em casos de emergência. A amplitude de aplicação dos problemas de localização de facilidades é realmente vasta. 

Tratando-se de vantagem competitiva no setor privado, conforme \citeonline{oliveira2012sistema}, “os problemas de localização de facilidades referem-se à decisão de onde posicionar os recursos que uma organização possui, de modo a atender, da melhor forma possível, aos critérios estabelecidos pela organização referentes aos pontos de demanda”. Esses critérios podem ser entendidos como restrições, que podem ser geográficas, operacionais, financeiras, de transporte, entre outras.

O problema de localização de facilidades está relacionado a diferentes tipos de cadeias de suprimentos, que podem variar desde uma configuração simples com um único CD, como discutido por \citeonline{SINGH2018343} no contexto da abertura de um armazém de peças automotivas no Irã, até cadeias mais complexas que envolvem fornecedores, CDs e PAs. Essas cadeias podem ser de circuito aberto, em que o produto segue do CD ao cliente, ou de circuito fechado, onde o produto vai do fornecedor ao cliente e pode retornar para recuperação ou descarte. O estudo \citeonline{AMIN2017316}, por exemplo, propõe um modelo de otimização multi-objetivo para a configuração de uma \sigla{CSCF}{Cadeia de Suprimentos de Circuito Fechado} global, incorporando os fluxos diretos e reversos.

Atualmente, observa-se um crescimento no consumo de produtos eletroeletrônicos. Conforme \citeonline{mindesvendaeletro}, em 2024, houve um aumento de 29\% na venda deste tipo de aparelho com um resultado de 117 milhões de unidades vendidas, sendo considerado o melhor desempenho na ultima década. Com esse tipo de expansão, surge a necessidade de gerenciar os resíduos decorrentes do descarte de materiais eletroeletrônicos. Atualmente, temos os \sigla{REEE}{Resíduos de Equipamentos Eletroeletrônicos}. Segundo \citeonline{Santos2014a}, essa abordagem tem recebido cada vez mais atenção da sociedade, pois esses resíduos contêm materiais que podem prejudicar tanto os seres humanos quanto o meio ambiente.

Existem alguns estudos que abordam a CSCF com LR voltada ao tratamento de REEE. Em seu estudo, \citeonline{demajorovic2016logistica} discutem os desafios e oportunidades relacionados à implementação da LR para computadores e celulares em países em desenvolvimento, como Índia, China e Brasil. Apesar do crescente interesse pelo tema nesses países, impulsionado pela necessidade de destinação adequada dos REEE, eles ainda se encontram em desvantagem em relação aos países desenvolvidos, principalmente devido à baixa conscientização dos consumidores.

Nesta dissertação, a pesquisa abordará o CSCF, destacando uma característica muito importante desse tipo de cadeia: a \sigla{LR}{Logística Reversa}. Conforme mencionado por \citeonline{guarnieri2011logistica}, a logística reversa é o processo que opera o retorno de produtos ao ciclo produtivo do negócio.A LR possui importância crucial conforme destacado por \citeonline{mogalelogirever}, pois cria valor na economia circular, ou seja, o produto não perde seu valor ao final da cadeia. Existe a possibilidade de receber o produto de volta e manter seu valor remanufaturando-o ou enviando para descartes e recuperando um percentual de valor através dos materiais reutilizeis em reciclagem.

%A \sigla{PLIM}{Programação Linear Inteira Mista} é uma abordagem amplamente utilizada na \sigla{PO}{Pesquisa Operacional} para a resolução de problemas complexos e variados. O estudo de \citeonline{mapa2012uso} demonstra que a aplicação de modelos de localização e alocação de facilidades pode proporcionar uma melhoria de até 37\% na minimização de custos, evidenciando que esse tipo de abordagem é essencial para o tipo de problema que será analisado neste estudo.

%A PLIM possui três componentes essenciais para resolver esse tipo de problema:

%\begin{itemize}
%    \item Função objetivo:  representa a meta a ser alcançada em um problema de otimização. Ela pode ser formulada para maximizar ou minimizar uma determinada quantidade, como lucro, custo, etc. A função objetivo depende das variáveis de decisão do problema, em que seus valores devem ser otimizados; 
%    \item Variáveis de decisão: representam as escolhas que podem ser tomadas considerando o espaço de soluções, e seus valores são otimizados para maximizar ou minimizar a função objetivo.;
%    \item Restrições: são as condições que o problema deve satisfazer, elas limitam as escolhas das variáveis de decisão para a resolução do problema.
%\end{itemize}

%Os três elementos fundamentais da programação matemática são aplicáveis à resolução de problemas de localização de facilidades. Para abordar essa resolução, é necessário primeiro definir a função objetivo, que pode ser tanto maximizar como minimizar alguma métrica, como a área a ser atendida ou os lucros obtidos. Em seguida, é importante identificar as variáveis de decisão, tais como o volume a ser movimentado e a abertura de novas facilidades, entre outras. Por fim, as restrições são determinadas pelo problema, incluindo capacidade operacional, área geográfica a ser atendida e outros aspectos relevantes.

Esta dissertação apresenta o desenvolvimento de um modelo matemático de PLIM aplicado à otimização de uma CSCF no setor de adquirência. O objetivo central do trabalho é a minimização do custo logístico operacional total em um horizonte de planejamento de um ano. Para isso, serão definidos quantos CDs e PAs serão abertos, qual será a localização geográfica e qual será a movimentação.

No modelo PLIM, busca-se a determinação de uma melhor solução factível para um conjunto de decisões estratégicas e operacionais cruciais. Essas decisões incluem a quantidade e as localizações geográficas dos CDs e PAs a serem ativados, bem como as quantidades de itens a serem movimentados e os níveis de estoque e armazenagem ao longo de toda a cadeia logística. A função objetivo abrange os custos de abertura, manutenção e movimentação (transporte) associados à infraestrutura (CDs e PAs).

No âmbito metodológico, o trabalho enfrentou desafios relacionados à capacidade computacional, dada a escala dos dados analisados para a programação do modelo de otimização. Para mitigar esta dificuldade e estruturar os dados com maior aderência à realidade, foram gerados três cenários, reduzidas em volume de dados, de simulação distintos, os quais incorporam variações de ruído na demanda do sistema.

A estrutura desta dissertação segue o seguinte formato: no Capítulo \ref{chapter:rev_bib}, apresenta-se uma revisão bibliográfica que explora temas como logística, LR e  problemas de localização de facilidades. No Capítulo \ref{chapter:def_prob_mod}, o problema é detalhado juntamente com o modelo matemático a ser otimizado. No Capítulo \ref{chapter:exp_comp} detalha os dados utilizados para gerar as instâncias utilizadas e também apresenta os resultados obtidos. Finalmente, as conclusões e trabalhos futuros estão no Capítulo \ref{chapter:conc}.
